\centerline{\bf \Huge Комбо база}
\centerline{\bf \Huge }

\begin{flushright}
        \scriptsize{Кто не будет решать тот лох!}
\end{flushright}

\problem{1}
В корзине лежат 12 яблок и 10 апельсинов. Ваня выбирает из нее либо яблоко, либо апельсин, после чего Надя берет и яблоко, и апельсин. В каком случае Надя имеет большую свободу выбора: если Ваня взял яблоко или если он взял апельсин?

\problem{2}
Сколько существует 5 -значных натуральных чисел, делящихся на 2 и не содержащих в десятичной записи ни одной из цифр $3,4,5,7$ и $9 ?$

\problem{3} \textbf{В этой задаче можно получить 1 балл, если решить пункты а и б (если решить только один из пунктов ничего получите) и 1 балл, если решить пункт в. То есть максимум 2 балла.}

a) Сколько существует 6 -значных чисел, в которых любые две цифры различны?

б) Сколько существует 6 -значных чисел, в которых есть цифра 7 , но нет цифры 9?

в)$^*$ Сколько существует 6 -значных чисел, у которых по три чётных и нечётных цифры?

\problem{4}
Назовем лесенкой высоты $n$ фигуру, состоящую из $n$ горизонтальных рядов: в нижнем ряду $2 n-1$ клетка, а в каждом следующем ряду слева и справа убирается по одной клетке. Сколькими способами на клетки лесенки высоты 10 можно поставить 10 ладей так, чтобы они не били друг друга?

\problem{5}
Андрей, Боря, Вася, Гриша, Денис и Женя после олимпиады собрались в кинотеатр. Они купили билеты на 6 мест подряд в одном ряду. Андрей и Боря хотят сидеть рядом, а Вася и Гриша не хотят. Сколькими способами они могут сесть на свои места с учётом их пожеланий?
